\documentclass[12pt,a4paper]{article}
\usepackage[utf8]{inputenc}
\usepackage[T1]{fontenc}
\usepackage[brazil]{babel}
\usepackage{geometry}
\usepackage{setspace}
\geometry{margin=2.5cm}
\setstretch{1.15}
\setlength{\parindent}{0pt}
\setlength{\parskip}{0.35em}

\newcommand{\resposta}[1]{\textbf{Resposta:} #1}

\begin{document}

\begin{center}
\textbf{Ministério da Educação}\\
\textbf{Universidade Tecnológica Federal do Paraná -- Campus Pato Branco}\\
\textbf{Departamento Acadêmico de Informática}\\[0.5em]
\textbf{Disciplina: Redes de Computadores} \hfill \textbf{Prof. Dr. Eden Dosciatti}

\vspace{1.2em}
\Large\textbf{Atividade Avaliativa 1}
\end{center}

\vspace{1em}
\noindent\textbf{Nome da(o) aluna(o):} Pedro Mariano

\vspace{1.2em}
\section*{TAREFA 1}
No curso de CCNA:Redes de Computadores, no NetAcad, no módulo 1, \textbf{1.0.5 Packet Tracer - Exploração de modo lógico e físico}, abra o arquivo em \textbf{Exploração de modo lógico e físico} e responda, aqui neste documento, as questões abaixo:

\subsection*{Parte 1}
\noindent\textbf{Quais são as subcategorias para dispositivos de rede?}

\resposta{Roteadores, switches, hubs, dispositivos sem fio, segurança e emulação WAN.}

\subsection*{Parte 2}
\noindent\textbf{Quais dispositivos usam uma conexão com fio para se conectar para alternar ALS2?}

\resposta{ALS1, Access\_Point e WebServer.}

\vspace{0.8em}
\noindent\textbf{Qual dispositivo está conectado ao Access\_Point?}

\resposta{Laptop\_1.}

\vspace{0.8em}
\noindent\textbf{Onde o dispositivo está conectado ao Access\_Point fisicamente localizado?}

\resposta{Na mesa (\textit{On the Table}).}

\vspace{1.2em}
\section*{TAREFA 2}
No curso de CCNA:Redes de Computadores, no NetAcad, no módulo 1, \textbf{1.5.5 Packet Tracer - Representação da Rede}, abra o arquivo em \textbf{Representação da Rede} e responda, aqui neste documento, as questões abaixo:

\subsection*{Etapa 1}
\noindent\textbf{Liste as categorias intermediárias de dispositivo.}

\resposta{Roteadores, switches, hubs, dispositivos sem fio e emulação WAN.}

\vspace{0.8em}
\noindent\textbf{Sem entrar na nuvem da Internet ou na intranet, quantos ícones na topologia representam dispositivos de terminal (apenas uma conexão levando a eles)?}

\resposta{15.}

\vspace{0.8em}
\noindent\textbf{Sem contar as duas nuvens, quantos ícones na topologia representam dispositivos intermediários (várias conexões que levam a eles)?}

\resposta{11.}

\vspace{0.8em}
\noindent\textbf{Quantos dispositivos finais não são computadores de mesa?}

\resposta{8.}

\vspace{0.8em}
\noindent\textbf{Quantos tipos diferentes de conexões de meio físico são usados nesta topologia de rede?}

\resposta{4.}

\subsection*{Etapa 2}
\noindent\textbf{Com base em seus estudos até agora, explique o modelo cliente-servidor.}

\resposta{Em redes modernas, um host pode atuar como cliente, servidor ou ambos, dependendo do software instalado. Servidores oferecem serviços e informações (como e-mail e páginas web), enquanto clientes solicitam e exibem essas informações. Um cliente pode se tornar servidor ao instalar software de servidor.}

\vspace{0.8em}
\noindent\textbf{Liste pelo menos duas funções de dispositivos intermediários.}

\resposta{Regenerar e retransmitir sinais de dados; manter informações sobre caminhos na rede; notificar erros/falhas de comunicação; redirecionar tráfego por caminhos alternativos quando há falha de enlace; classificar e encaminhar mensagens por prioridades de QoS; permitir ou negar fluxo de dados conforme políticas de segurança.}

\vspace{0.8em}
\noindent\textbf{Liste pelo menos dois critérios para escolher um tipo de meio físico de rede.}

\resposta{Distância máxima de transmissão do sinal; ambiente de instalação; quantidade de dados e velocidade necessária; custo do meio e da instalação.}

\subsection*{Etapa 3}
\noindent\textbf{Explique a diferença entre uma LAN e uma WAN. Dê exemplos de cada uma.}

\resposta{LAN conecta usuários em área geográfica pequena (ex.: casa, escritório, campus). WAN conecta redes/usuários em áreas amplas e longas distâncias (ex.: MAN, Internet e intranets corporativas entre filiais).}

\vspace{0.8em}
\noindent\textbf{Na rede do Packet Tracer, quantas WANs você vê?}

\resposta{2 (Internet e Intranet).}

\vspace{0.8em}
\noindent\textbf{Quantas LANs você vê?}

\resposta{3.}

\vspace{0.8em}
\noindent\textbf{A Internet nesta rede Packet Tracer é excessivamente simplificada e não representa a estrutura e a forma da internet real. Descreva brevemente a Internet.}

\resposta{A Internet é uma malha global de redes interconectadas (internetworks), usada para comunicação com recursos que estão em outras redes.}

\vspace{0.8em}
\noindent\textbf{Quais são algumas das maneiras comuns de um usuário doméstico se conectar à Internet?}

\resposta{Cabo, DSL, discada, celular e satélite.}

\vspace{0.8em}
\noindent\textbf{Quais são alguns métodos comuns que as empresas usam para se conectar à Internet em sua área?}

\resposta{Link dedicado (leased line), Metro-Ethernet, DSL, cabo e satélite.}

\end{document}
